\section{Timeline Layout Engine}
\label{sec:timeline-rendering}

This section specifies the layout engine for the Timeline Grammar---the deterministic transformation from Timeline Domain IR to Scene IR. The layout engine is grammar-specific; the rendering backends (Section~\ref{sec:backends}) are shared across all grammars.

\subsection{Layout Pipeline Overview}

The Timeline layout engine defines a \emph{deterministic mapping} from a Timeline IR document and a
theme to Scene IR geometry: positions, dimensions, label placements, and visual
treatments. ``Deterministic'' is not aspirational; it is a hard contract. Two conforming
layout engines, given identical IR and theme, must produce geometrically identical Scene IR~\cite{vegalite2017}.

The engine is organized as an ordered six-phase pipeline. Each phase consumes only outputs
of preceding phases and produces geometry for later phases. The pipeline is \emph{pure}: no
phase reads from the IR after its designated turn, and no phase performs I/O, reads a system
clock, or invokes a random number generator.

\subsection{Determinism Contract}
\label{sec:rendering-determinism}

Two conforming layout engines $R_1$ and $R_2$, given the same IR document $D$ and theme $\theta$,
must produce identical coordinate values, label positions, and visual-treatment assignments.
This contract requires:

\begin{enumerate}
  \item \textbf{Pure function.}
        Output is a pure function of $(D,\,\theta)$. No external state---timestamps, random
        values, environment variables, system locale, or screen resolution---influences any
        computed value.

  \item \textbf{Stable sort keys.}
        Every ordered collection is sorted by an explicit, unambiguous key sequence.
        Collections without semantic ordering sort by \texttt{id} alphabetically:
        \begin{itemize}
          \item Tracks: by \texttt{index} ascending (integer, unique by well-formedness invariant).
          \item Activities within a track: by \texttt{(start\_ordinal,\;id)} ascending.
          \item Milestones: by \texttt{(date\_ordinal,\;id)} ascending.
          \item Annotations, sections, groups: by \texttt{id} ascending.
        \end{itemize}
        The ID values break all ties; global uniqueness is guaranteed by the IR invariant.

  \item \textbf{Fixed rounding.}
        All intermediate floating-point values are rounded using \emph{round-half-up}
        ($\lfloor v + 0.5 \rfloor$). Scene geometry values are expressed to two decimal
        places using the same convention. No other rounding rule is used anywhere in the
        pipeline.

  \item \textbf{Deterministic symbolic date resolution.}
        The symbolic date \texttt{now} and relative dates (\texttt{+3m}, \texttt{-2w})
        resolve to a fixed anchor. The anchor is: \texttt{metadata.today} if present;
        else \texttt{metadata.created}; else a validation error (see Gap~1 in
        \S\ref{sec:rendering-gaps}). The system clock is never consulted.

  \item \textbf{Embedded font metrics.}
        Label-width and line-height computations use metrics from font files bundled with
        the theme. System fonts are not consulted for any layout-critical measurement.
        This guarantees that label placement is identical on macOS, Linux, and Windows.

  \item \textbf{Specification-version governance.}
        The renderer reads \texttt{version} from the IR document and verifies it against a
        supported specification version before beginning layout. A version mismatch is a
        hard error, not a warning.

  \item \textbf{Layered determinism.}
        The determinism contract applies at three levels. (i)~\emph{Scene geometry} is
        always byte-deterministic: the layout pipeline output is bitwise-reproducible
        across runs, platforms, and renderer implementations. (ii)~\emph{Per-backend
        output} is deterministic given a pinned backend version and fixed effect seeds;
        the backend version is an explicit parameter of the contract.
        (iii)~\emph{Cross-backend pixel identity} is explicitly \emph{not} promised: the
        SVG backend and the Raster backend are permitted and expected to produce visually
        different outputs from the same Scene. This is correct behaviour at different
        fidelity tiers, not a defect. See Section~\ref{sec:scene-determinism} for the
        full layered contract.
\end{enumerate}

\paragraph{Consequence.}
If $R_1$ and $R_2$ both satisfy this contract and their outputs diverge, the divergent
renderer has a defect. The specification is the arbiter.

\subsection{Coordinate System}
\label{sec:rendering-coords}

The renderer works in a \emph{logical-pixel} coordinate system. The top-left corner of
the canvas is the origin $(0,\,0)$; $x$ increases rightward, $y$ downward. All layout
computations use integer arithmetic in logical pixels; the Scene/Render IR records
decimal-valued coordinates as its output (see Section~\ref{sec:scene-output}).

\begin{figure}[h]
\centering
\begin{Verbatim}[frame=single,fontsize=\small]
  x=0           x=Hhdr+mL              x=canvas.width
  +--------------+----------------------------------+
  |              |  Axis header (Haxis px)          |  y = mT
  |  Track       +----------------------------------+
  |  header      |  Track[0]  (row_height px)       |
  |  column      |  [activity bars, milestones]     |
  |  (Hhdr px)   +----------------------------------+
  |              |  Track[1]  (row_height px)       |
  |              |    ...                            |
  |              +----------------------------------+
  |              |  Track[N-1]                      |
  +--------------+----------------------------------+
                                         y = canvas.height
\end{Verbatim}
\caption{Canvas coordinate system and layout zones}
\label{fig:canvas-layout}
\end{figure}

\begin{table}[h]
\centering
\small
\begin{tabular}{lp{9.5cm}}
\toprule
\textbf{Symbol} & \textbf{Definition} \\
\midrule
$W$ & \texttt{theme.canvas.width} --- total canvas width in logical pixels \\
$m_L,\,m_R,\,m_T,\,m_B$ & Left, right, top, bottom margins \\
$H_{\mathrm{hdr}}$ & \texttt{theme.track.header\_width} --- left column for track labels \\
$H_{\mathrm{axis}}$ & \texttt{theme.axis.height} --- horizontal time-axis strip \\
$W_{\mathrm{draw}}$ & $W - m_L - m_R - H_{\mathrm{hdr}}$ --- drawable timeline width \\
$H_{\mathrm{draw}}$ & $\displaystyle\sum_{i} h_i + \sum_{i>0} g$ --- total height of all track rows plus gaps \\
$H$ & Computed canvas height: $m_T + H_{\mathrm{axis}} + H_{\mathrm{draw}} + m_B$ \\
\bottomrule
\end{tabular}
\caption{Coordinate system symbols}
\label{tab:coord-symbols}
\end{table}

The canvas \emph{width} is fixed by the theme. The canvas \emph{height} is always computed
from track count and sub-lane expansion; it must never be truncated.

\subsection{Timeline Axis}
\label{sec:rendering-axis}

\subsubsection{Axis Unit and Inference}
\label{sec:axis-inference}

\texttt{metadata.axis\_unit} controls tick granularity. If absent, the renderer infers a
suitable unit from the \texttt{time\_range} span, applying the first matching threshold:

\begin{table}[h]
\centering
\begin{tabular}{ll}
\toprule
\textbf{time\_range span (days, inclusive)} & \textbf{Inferred axis\_unit} \\
\midrule
$\leq 14$  & \texttt{day}     \\
$\leq 91$  & \texttt{week}    \\
$\leq 548$ & \texttt{month}   \\
$\leq 1461$& \texttt{quarter} \\
$\leq 2922$& \texttt{half}    \\
$> 2922$   & \texttt{year}    \\
\bottomrule
\end{tabular}
\caption{Axis unit inference thresholds (first match wins)}
\label{tab:axis-inference}
\end{table}

These thresholds yield 4--26 visible ticks at a 1200-pixel canvas width---the range that
avoids crowding while remaining informative for the executive-roadmap use cases described
in Section~\ref{sec:principles}.

\subsubsection{Date-to-Coordinate Function}
\label{sec:date-to-x}

All horizontal positions derive from a single deterministic function $x(T)$. Let
$T_{\mathrm{ord}}$ denote the \emph{day ordinal} of date $T$: an integer count of days
since 2000-01-01 (epoch day~0). Coarser-precision dates are resolved to their period-start
day before ordinal conversion (see \S\ref{sec:date-coercion}).

\[
  x(T) \;=\; H_{\mathrm{hdr}} + m_L \;+\;
  \left\lfloor
    \dfrac{
      (T_{\mathrm{ord}} - T_{s,\mathrm{ord}}) \;\cdot\; W_{\mathrm{draw}}
    }{
      T_{e,\mathrm{ord}} - T_{s,\mathrm{ord}}
    }
    \;+\; 0.5
  \right\rfloor
\]

where $T_{s,\mathrm{ord}}$ and $T_{e,\mathrm{ord}}$ are the day ordinals of the axis domain
start and end after coercion. The numerator product
$(T_{\mathrm{ord}} - T_{s,\mathrm{ord}}) \cdot W_{\mathrm{draw}}$ is computed before division
to prevent precision loss. For a 1200-pixel canvas spanning ten years (3\,650 days), the
intermediate product is at most $1200 \times 3650 = 4\,380\,000$, well within 32-bit
integer range.

\paragraph{Boundary invariants.}
$x(T_s) = H_{\mathrm{hdr}} + m_L$ exactly;
$x(T_e) = H_{\mathrm{hdr}} + m_L + W_{\mathrm{draw}}$ exactly. Any date at or before $T_s$
maps to the left canvas boundary; any date at or after $T_e$ maps to the right canvas
boundary, before clipping.

\subsubsection{Date Coercion for Bar Edges}
\label{sec:date-coercion}

Dates coarser than \texttt{day} must be coerced to a specific calendar day before ordinal
conversion. The rule depends on whether the date is a \emph{left edge} (start of a bar) or
a \emph{right edge} (end of a bar):

\begin{table}[h]
\centering
\small
\begin{tabular}{lll}
\toprule
\textbf{Date form} & \textbf{As left edge (start)} & \textbf{As right edge (end)} \\
\midrule
\texttt{2026-06-09}  (day)    & 2026-06-09 & 2026-06-09 \\
\texttt{2026-06}     (month)  & 2026-06-01 & 2026-06-30 \\
\texttt{2026}        (year)   & 2026-01-01 & 2026-12-31 \\
\texttt{2026-Q2}     (quarter)& 2026-04-01 & 2026-06-30 \\
\texttt{2026-H1}     (half)   & 2026-01-01 & 2026-06-30 \\
\texttt{FY26-Q2}     (fiscal) & First day of fiscal Q2 & Last day of fiscal Q2 \\
\bottomrule
\end{tabular}
\caption{Date coercion to day boundaries by edge role}
\label{tab:date-coercion}
\end{table}

\textbf{Fiscal dates} require a fiscal-year start month that is currently absent from the IR
contract (see Gap~2 in \S\ref{sec:rendering-gaps}).

The same left-edge/right-edge coercion applies to \texttt{time\_range.start} (left) and
\texttt{time\_range.end} (right) when establishing the axis domain.

\subsubsection{Tick and Gridline Placement}
\label{sec:tick-placement}

\begin{enumerate}
  \item Enumerate all \texttt{axis\_unit} period-start dates within $[T_s,\,T_e]$ inclusive.
  \item For each period boundary $T_k$: compute $x_k = x(T_k)$ per \S\ref{sec:date-to-x}.
  \item Place a tick mark of height \texttt{theme.axis.tick\_height}~px at $x_k$, at the
        bottom of the axis header band.
  \item Place a vertical gridline at $x_k$ spanning from the top of Track[0] to the bottom
        of Track[$N-1$], styled per \texttt{theme.axis.gridline\_*} properties.
  \item Place a tick label centered at $x_k$, \texttt{theme.axis.tick\_label\_offset}~px
        below the tick. Format using \texttt{metadata.locale} and the label rules in
        Table~\ref{tab:tick-label-density}. If a label would overlap its neighbor (right
        edge $>$ neighbor's left edge), suppress the label for this tick only.
\end{enumerate}

\begin{table}[h]
\centering
\small
\begin{tabular}{lll}
\toprule
\textbf{axis\_unit} & \textbf{Primary label content} & \textbf{Secondary label (every $N$ ticks)} \\
\midrule
\texttt{day}     & Day-of-month number         & Month name (every 7) \\
\texttt{week}    & Week-start day              & Month name (every 4) \\
\texttt{month}   & Month abbreviation          & Year number (every 12) \\
\texttt{quarter} & Q1 / Q2 / Q3 / Q4           & Year number (every 4) \\
\texttt{half}    & H1 / H2                     & Year number (every 2) \\
\texttt{year}    & Year number                 & None \\
\bottomrule
\end{tabular}
\caption{Tick label content rules per axis unit}
\label{tab:tick-label-density}
\end{table}

\paragraph{Week ticks.}
For \texttt{axis\_unit = week}, week boundaries are ISO~8601 Monday starts regardless of
locale. The locale affects label string formatting only (language, order), never tick positions.

\paragraph{Section bands.}
If \texttt{sections[]} is non-empty, the theme may render alternating column shading between
section boundaries. Band edges coincide with section start/end dates coerced per
Table~\ref{tab:date-coercion}. Section bands are painted at the lowest z-order, below
gridlines and all content elements.

\subsection{The Six-Phase Layout Algorithm}
\label{sec:layout-algorithm}

The renderer executes geometry computation in exactly six phases in the order given below.
Each phase's output is immutable input to subsequent phases. No phase may invoke logic from
a later phase.

\begin{lstlisting}[caption={Layout pipeline overview}]
Phase 1 | AXIS COMPUTATION
        |   Input:  IR metadata (time_range, axis_unit, today-anchor)
        |   Output: axis domain [Ts, Te], axis_unit, tick_positions[]

Phase 2 | TRACK PLACEMENT
        |   Input:  tracks[] sorted by index
        |   Output: track.y_top[], track.height[] (provisional)

Phase 3 | ACTIVITY GEOMETRY
        |   Input:  activities[], Phase-1 axis, Phase-2 track positions
        |   Output: activity.{x_left, x_right, y, lane, width}
        |           track.height[] (final, expanded for sub-lanes)

Phase 4 | MILESTONE GEOMETRY
        |   Input:  milestones[], Phase-1 axis, Phase-3 final track heights
        |   Output: milestone.{x_center, y_center, stack_index}

Phase 5 | SECTIONS AND ANNOTATIONS
        |   Input:  sections[], annotations[], all Phase 1-4 geometry
        |   Output: section.{x_left, x_right}, annotation.{x, y}

Phase 6 | LABEL COLLISION RESOLUTION
        |   Input:  all geometry from Phases 3-5
        |   Output: final label positions for milestones
\end{lstlisting}

\subsubsection{Phase 1: Axis Computation}
\label{sec:phase1}

\begin{lstlisting}[caption={Phase 1: Axis computation}]
1. Resolve today-anchor:
     if metadata.today is set : anchor = metadata.today
     elif metadata.created is set : anchor = metadata.created
     else : ERROR("Cannot resolve symbolic/relative dates: no today/created")

2. Resolve time_range.start:
     a. Expand 'now' -> anchor; expand '+3m' -> anchor + 3 months; etc.
     b. Coerce to period start (left-edge rule, Table 5.3).  -> T_s

3. Resolve time_range.end (same expansion; coerce to period end -> T_e).
     ASSERT T_e >= T_s; else ERROR.

4. Infer axis_unit if absent:
     span_days = T_e_ord - T_s_ord
     Apply Table 5.2 thresholds in order; assign axis_unit.

5. Compute W_draw = canvas.width - mL - mR - theme.track.header_width.

6. Enumerate tick positions:
     T_k = all period-start dates of axis_unit in [T_s, T_e] inclusive.
     For each T_k: x_k = x(T_k) per Section 5.3.2.
     Compute tick label strings using metadata.locale.
\end{lstlisting}

\subsubsection{Phase 2: Track Placement}
\label{sec:phase2}

\begin{lstlisting}[caption={Phase 2: Track placement}]
1. Sort tracks by index ascending.
   ASSERT all index values are unique non-negative integers.

2. y_cursor = mT + Haxis

3. For i in 0 .. N_tracks-1:
     track[i].y_top  = y_cursor
     track[i].height = theme.track.row_height     -- provisional
     y_cursor       += theme.track.row_height + theme.track.row_gap

4. Record group bracket extents after Phase 3 finalises track heights.
\end{lstlisting}

Track heights set here are provisional. Phase~3 may expand them when overlapping activities
require sub-lanes. Downstream phases must use the Phase-3-finalised heights.

\subsubsection{Phase 3: Activity Geometry}
\label{sec:phase3}

\paragraph{Sub-lane assignment.}

\begin{lstlisting}[caption={Phase 3: Activity geometry and greedy sub-lane assignment}]
For each track T (in index order):
  1. A = activities where A.track == T.id.
  2. Resolve each A.start_ord = ordinal(coerce_left(A.start)).
     Resolve each A.end_x per special-end-date rules below.
  3. Sort A by (start_ord, id) ascending (stable sort).
  4. Greedy sub-lane assignment:
       lane_end_x = [-infinity]
       For each A_i in sort order:
         k = smallest index where lane_end_x[k] <= x(A_i.start_ord)
         if no such k exists:
           if len(lane_end_x) < theme.track.max_sub_lanes:
             append -infinity to lane_end_x; k = new last index
           else:
             k = theme.track.max_sub_lanes - 1   -- cap; emit WARNING
         A_i.lane    = k
         lane_end_x[k] = A_i.end_x
  5. n_lanes = max(A.lane) + 1
  6. if n_lanes > 1:
       track[T].height = theme.track.row_height
                       + (n_lanes - 1) * theme.track.sub_lane_height
  7. For each A_i:
       A_i.y = track[T].y_top
             + theme.activity.bar_top_margin
             + A_i.lane * theme.track.sub_lane_height
       A_i.x_left  = x(A_i.start_ord)
       A_i.x_right = A_i.end_x
       logical_w   = A_i.x_right - A_i.x_left
       if logical_w < theme.activity.min_width:
         mid         = (A_i.x_left + A_i.x_right) / 2   -- round-half-up
         A_i.x_left  = mid - theme.activity.min_width / 2
         A_i.x_right = A_i.x_left + theme.activity.min_width
       A_i.width = A_i.x_right - A_i.x_left
\end{lstlisting}

\paragraph{Special end-date resolution.}
\begin{itemize}
  \item \texttt{end} absent (omitted): treated identically to \texttt{ongoing}.
  \item \texttt{ongoing}: \texttt{end\_x} $=$ $H_{\mathrm{hdr}} + m_L + W_{\mathrm{draw}}$
        (right canvas boundary); append right-chevron ($\rightarrow$) decoration.
  \item \texttt{tbd}: \texttt{end\_x} $=$ \texttt{x\_left} $+$
        \texttt{theme.activity.tbd\_extension\_px}
        (default: $2 \times \overline{\Delta x_{\text{tick}}}$, the mean tick spacing in pixels);
        render rightward extension with dashed border at 50\% opacity;
        place ``TBD'' label inside extension zone.
  \item \texttt{unknown}: same geometry as \texttt{tbd}; use \texttt{unknown} visual
        treatment from theme status map.
  \item \texttt{\textasciitilde{}date}: coerce date normally to ordinal; \texttt{end\_x}
        $= x(\text{date\_ord})$; flag right edge as approximate for gradient-fade rendering.
\end{itemize}

\paragraph{Label placement (deterministic three-way rule).}
The rule is applied exactly once per activity; no iteration:
\begin{enumerate}
  \item \textbf{Inside}: if \texttt{A.width} $\geq$ \texttt{theme.activity.label\_inside\_min\_width},
        render label horizontally centered in bar, vertically centered on bar height, in
        \texttt{theme.activity.label\_color\_inside}.
  \item \textbf{Right}: render label starting at $\texttt{A.x\_right} + 4$~px in
        \texttt{theme.activity.label\_color\_outside}.
        If label right edge $>$ $H_{\mathrm{hdr}} + m_L + W_{\mathrm{draw}} - 4$, proceed to
        step~3.
  \item \textbf{Left}: render label ending at $\texttt{A.x\_left} - 4$~px.
        If label left edge $<$ $H_{\mathrm{hdr}} + m_L$, truncate label to
        \texttt{theme.activity.label\_truncate\_chars} characters $+$
        ``\ldots'' (U+2026) and re-attempt step~2.
\end{enumerate}

\paragraph{Progress indicator.}
If \texttt{activity.progress} is set: render a filled strip of height
\texttt{theme.activity.progress\_bar\_height}~px at the bottom of the bar, width $=
\lfloor \texttt{A.width} \cdot \texttt{A.progress} + 0.5 \rfloor$~px, inset inside the
bar border.

\subsubsection{Phase 4: Milestone Geometry}
\label{sec:phase4}

\paragraph{Placement and stacking.}

\begin{lstlisting}[caption={Phase 4: Milestone placement and stacking}]
For each milestone M (sorted by date_ord, id):
  1. x_center = x(M.date_ordinal).
  2. Determine base_y:
       if M.track is set:
         base_y = track[M.track].y_top + track[M.track].height / 2
       else (cross-track):
         base_y = mT + Haxis + H_draw / 2

  3. Stacking within (effective_track_id, x_center) group:
       Sort group members by id ascending.
       Assign stack_index k = 0, 1, 2, ...
       M.y_center = base_y + k * theme.milestone.stack_offset_y

  4. Diamond vertices (integer arithmetic):
       d = theme.milestone.diamond_size
       top    = (x_center,     M.y_center - d)
       right  = (x_center + d, M.y_center    )
       bottom = (x_center,     M.y_center + d)
       left   = (x_center - d, M.y_center    )
\end{lstlisting}

The diamond is a four-vertex polygon with vertices computed directly. No
\texttt{transform="rotate(...)"} attribute is used; rotation-transform-dependent rendering
differences across SVG viewers are avoided.

\paragraph{Cross-track milestones.}
When \texttt{track} is null, the milestone is drawn at the vertical centre of the full
rendered area. Its diamond size is \texttt{theme.milestone.diamond\_size} (not scaled);
visual emphasis comes from the spanning vertical extent of the label, which may be increased
by a theme multiplier \texttt{theme.milestone.cross\_track\_label\_scale} (default: 1.2).

\subsubsection{Phase 5: Sections and Annotations}
\label{sec:phase5}

\paragraph{Section bands.}
For each section $S$ (sorted by start date, then id):
\begin{itemize}
  \item $x_{\text{left}}(S) = x(\text{coerce\_left}(S.\texttt{start}))$.
  \item $x_{\text{right}}(S) = x(\text{coerce\_right}(S.\texttt{end}))$.
  \item Band spans vertically from $m_T + H_{\mathrm{axis}}$ to
        $m_T + H_{\mathrm{axis}} + H_{\mathrm{draw}}$.
  \item Z-order: sections paint first (below gridlines, activities, milestones, labels).
\end{itemize}
Overlapping sections are valid; a later section in sort order paints over an earlier one.

\paragraph{Annotation placement.}
\begin{enumerate}
  \item Sort annotations by \texttt{id} ascending.
  \item For each annotation: compute anchor as centre of target entity's bounding box.
  \item Place annotation box in the first quadrant that keeps it entirely within the canvas.
        Preference order: top-right, top-left, bottom-right, bottom-left.
        If no quadrant fits, use top-right and clip to canvas boundary.
  \item Draw connector (line/elbow) from annotation box edge to anchor,
        styled per \texttt{theme.annotation.connector\_style}.
\end{enumerate}

\subsubsection{Phase 6: Label Collision Resolution}
\label{sec:phase6}

Activity labels are placed definitively in Phase~3 (no iteration). Milestone labels require
a collision-resolution pass because multiple milestones at similar $x$ positions can produce
overlapping labels.

\begin{lstlisting}[caption={Phase 6: Milestone label collision resolution}]
1. Collect all milestone label bounding boxes (x_label, y_label, width, height, id).
2. Sort by (x_label, y_label, id) ascending (stable).
3. Scan pass (bounded by N iterations, N = label count):
     For each consecutive pair (L_i, L_{i+1}):
       if L_i.right_edge > L_{i+1}.left_edge
          AND |L_i.y_label - L_{i+1}.y_label| < label_height:
         L_{i+1}.y_label += theme.milestone.label_stack_offset
4. Repeat step 3 until no overlapping pairs remain OR N iterations reached.
   If overlaps persist after N iterations: truncate labels to
   theme.milestone.label_max_width px, accepting visual overlap.
5. Clamp all final y_label values to [mT, mT + Haxis + H_draw + mB].
\end{lstlisting}

Because the input sort and every shift decision are deterministic, two renderers executing
this algorithm on identical inputs produce identical label placements.

\subsection{Scene / Render IR --- Output of the Pipeline}
\label{sec:scene-output}

The six-phase layout pipeline does not target any output format directly. Its output is the
\textbf{Scene / Render IR}: a byte-deterministic, backend-agnostic record of every drawing
primitive, resolved coordinate, visual treatment, and effect request. The Scene is described
fully in Section~\ref{sec:scene-ir}; its role here is to name what the pipeline produces
and to record what each phase contributes.

Every value computed in Phases 1--6 becomes a field of a Scene drawing primitive or a
top-level Scene property. The Scene is a pure data record: it carries no rendering
instructions specific to SVG, rasterisation, or PowerPoint. Rendering backends (SVG,
Raster, PPTX, HTML) consume the Scene and emit their respective formats; they are described
in Section~\ref{sec:outputs}.

\paragraph{What the Scene preserves from the pipeline.}
\begin{itemize}
  \item \textbf{Phase 1} outputs: axis domain $[T_s,\,T_e]$, axis unit, tick positions
        $x_k$, and tick label strings.
  \item \textbf{Phases 2--3} outputs: all activity bar coordinates
        $(x_{\text{left}},\,x_{\text{right}},\,y,\,\text{width})$, lane assignments,
        and final track heights.
  \item \textbf{Phase 4} outputs: all milestone vertex coordinates and stack indices.
  \item \textbf{Phase 5} outputs: section band bounds
        $(x_{\text{left}},\,x_{\text{right}})$, annotation positions, and connector
        endpoints.
  \item \textbf{Phase 6} outputs: final label positions after collision resolution.
  \item \textbf{Theme-resolved visual treatments}: fill colours, stroke styles, opacity,
        corner radii, pattern references, and effect requests with their fallback policies
        (see Section~\ref{sec:schema-fidelity}).
\end{itemize}

The Scene does not carry the IR document or theme; it is the fully resolved,
self-contained rendering specification. A backend needs only the Scene and its own
capability profile to produce output.

\subsection{Edge Cases}
\label{sec:edge-cases}

Every case below is normative. A renderer that produces different behaviour for any case
has a defect. The summary table gives the complete catalogue; extended descriptions follow
for the most structurally complex cases.

\begin{longtable}{@{}p{3.2cm}p{3.4cm}p{5.8cm}@{}}
\caption{Edge-case definitions (normative)}
\label{tab:edge-cases} \\
\toprule
\textbf{Case} & \textbf{Condition} & \textbf{Rendering Rule} \\
\midrule
\endfirsthead
\multicolumn{3}{c}{\small\textit{Table~\ref{tab:edge-cases} continued}} \\[2pt]
\toprule
\textbf{Case} & \textbf{Condition} & \textbf{Rendering Rule} \\
\midrule
\endhead
\midrule
\multicolumn{3}{r}{\small\textit{continued \ldots}} \\
\endfoot
\bottomrule
\endlastfoot
Zero-duration activity &
  \texttt{start == end} or \texttt{span} that resolves to a single point &
  Render bar of \texttt{min\_width}~px centred at $x(\text{start})$. Never render 0-width. Label placed right (outside). \\[4pt]
Overlapping activities (same track) &
  Two activities in same track whose date ranges overlap &
  Greedy sub-lane assignment (Phase 3). Track height expands. Cap at \texttt{max\_sub\_lanes}. \\[4pt]
Open interval (\texttt{ongoing}) &
  \texttt{end: ongoing} or \texttt{end} absent &
  Bar extends to right canvas boundary. Right-chevron decoration. No overflow beyond canvas edge. \\[4pt]
TBD end date &
  \texttt{end: tbd} &
  Bar extends \texttt{tbd\_extension\_px} beyond \texttt{x\_left}. Dashed, 50\% opacity. ``TBD'' label inside extension. \\[4pt]
Both dates TBD &
  \texttt{start: tbd} or equivalent &
  Full-track-width hatched block at 40\% opacity. Activity label $+$ ``(TBD)''. \\[4pt]
Unknown start &
  \texttt{start: unknown} &
  Left edge placed at $x(T_s)$. Left-fade gradient over \texttt{approx\_fade\_px}. \\[4pt]
Unknown end &
  \texttt{end: unknown} &
  Same geometry as \texttt{tbd}; rendered with \texttt{unknown} visual treatment from theme. \\[4pt]
Approximate start &
  \texttt{start: \textasciitilde{}date} &
  Geometry at nominal date; left edge rendered with gradient fade of \texttt{approx\_fade\_px}. \\[4pt]
Approximate end &
  \texttt{end: \textasciitilde{}date} &
  Geometry at nominal date; right edge rendered with gradient fade. \\[4pt]
Activity fully before range &
  $\text{end} < T_s$ &
  Not rendered. Renderer warning. Activity still appears in legend. \\[4pt]
Activity fully after range &
  $\text{start} > T_e$ &
  Not rendered. Renderer warning. Activity still appears in legend. \\[4pt]
Activity partially before range &
  $\text{start} < T_s$, $\text{end} \geq T_s$ &
  Bar clipped to $[x(T_s),\;x(\text{end})]$. Left-clip indicator on clipped edge. \\[4pt]
Activity partially after range &
  $\text{start} \leq T_e$, $\text{end} > T_e$ &
  Bar clipped to $[x(\text{start}),\;x(T_e)]$. Right-clip indicator on clipped edge. \\[4pt]
Very short span &
  $x(\text{end}) - x(\text{start}) < \text{min\_width}$ after rounding &
  Render at \texttt{min\_width}; centre on logical midpoint (round-half-up). Label right (outside). \\[4pt]
Very long span &
  Bar occupies $> 80\%$ of $W_{\mathrm{draw}}$ &
  Render normally. Prefer inside label placement if label fits within bar. \\[4pt]
Simultaneous milestones (same track) &
  Two milestones with equal \texttt{date} on the same \texttt{track} &
  Stack vertically by \texttt{stack\_offset\_y}, sorted by \texttt{id} ascending (Phase 4). \\[4pt]
Simultaneous milestones (cross-track) &
  Two milestones with \texttt{track: null} and equal \texttt{date} &
  Stack at full-timeline vertical centre, sorted by \texttt{id} ascending. \\[4pt]
Milestone outside range &
  Milestone date $< T_s$ or $> T_e$ &
  Not rendered. Renderer warning. Appears in legend. \\[4pt]
Empty track &
  Track has no activities and no milestones &
  Row rendered at \texttt{row\_height}. Label visible. Body empty. Never collapsed. \\[4pt]
Sub-lane cap exceeded &
  Overlapping activities require $>$ \texttt{max\_sub\_lanes} lanes &
  Excess activities placed in the last lane (visible overlap). Renderer warning. \\
\end{longtable}

\paragraph{Overlapping activities: sub-lane assignment in detail.}

The greedy algorithm assigns each activity (in stable \texttt{(start\_ord,\;id)} order) to the
lowest-indexed lane where it does not overlap any previously placed activity. Overlap is
defined as $x_{\text{left}}(A_j) < x_{\text{right}}(A_i)$ for activities $A_i$, $A_j$ in
the same lane. Because sort order is deterministic and each assignment is a pure function of
the current \texttt{lane\_end\_x} state, all renderers produce identical lane assignments
for identical inputs.

\paragraph{Simultaneous milestones in detail.}

When milestones share an $(\texttt{track},\,x_{\text{center}})$ pair, they are stacked downward
from the nominal track centre. The first milestone in \texttt{id}-ascending order is at the
base position (stack index 0). Subsequent milestones are offset by
$k \times \texttt{theme.milestone.stack\_offset\_y}$. Stacking is permitted to overflow into
the inter-track gap; milestones must \emph{not} be silently suppressed if they cannot fit
within track height.

\paragraph{Clip indicators.}

An activity bar clipped at the axis boundary receives a \emph{clip indicator}: a small
``/\!/'' mark (two angled lines, $4 \times 8$~px) drawn in the activity's fill colour at
full opacity, positioned at the clipped edge centre. The indicator signals to the reader
that the activity continues beyond the visible axis range.

\paragraph{Minimum-width semantics.}

\texttt{theme.activity.min\_width} (default: 4~px) prevents activities from being rendered
invisibly narrow at coarse zoom levels. A zero-duration activity and a very-short-span
activity both produce a bar of exactly \texttt{min\_width} px. The two cases are
\emph{visually indistinguishable}; distinguishing them would require a separate icon or
tooltip, which is a theme-level option outside the core rendering contract.

\subsection{Known IR Gaps}
\label{sec:rendering-gaps}

The following gaps in Mark's IR contract (\S\ref{sec:ir}) affect rendering determinism.
They are flagged here for Mark and Leslie to resolve; this spec does \emph{not} unilaterally
extend the IR. Renderers encountering these gaps should apply the documented fallback and
emit a validation warning.

\begin{description}
  \item[Gap 1 --- \texttt{metadata.today} field missing.]
    Leslie's decision record mandates an explicit \texttt{today} field in the IR for
    deterministic resolution of the \texttt{now} symbolic date and the today-marker
    feature. Mark's current \texttt{metadata} object does not include this field.
    \textbf{Proposed resolution:} add \texttt{today} (type \texttt{date}, optional) to
    \texttt{metadata}. Renderer fallback chain:
    \texttt{metadata.today} $\to$ \texttt{metadata.created} $\to$ validation error.

  \item[Gap 2 --- \texttt{metadata.fiscal\_year\_start} missing.]
    The \texttt{FY26-Q2} fiscal-date format requires knowing the fiscal-year start month.
    Two renderers assuming different organisations' fiscal calendars (US federal: October;
    UK: April; common corporate: January) will produce different $x$-coordinates for
    identical fiscal dates.
    \textbf{Proposed resolution:} add \texttt{fiscal\_year\_start} (type \texttt{int},
    range 1--12, default 1 = January) to \texttt{metadata}.

  \item[Gap 3 --- Relative date anchor undefined.]
    Relative dates (\texttt{+3m}, \texttt{-2w}) are described in Mark's contract as
    resolving from ``context'', which is insufficient. A renderer must know the anchor date
    to compute an ordinal.
    \textbf{Proposed resolution:} define that relative dates use the same anchor as Gap~1
    (\texttt{today} $\to$ \texttt{created} $\to$ error).

  \item[Gap 4 --- Omitted \texttt{end} semantics ambiguous.]
    Mark's Activity contract marks \texttt{end} as optional but does not state whether an
    absent \texttt{end} is equivalent to \texttt{ongoing} or a validation error. This
    rendering model treats omitted \texttt{end} as \texttt{ongoing} (see the open-interval
    row in Table~\ref{tab:edge-cases}). The IR spec should make this explicit.
\end{description}
